\documentclass[11pt]{article}

\usepackage[margin=1in]{geometry}
\usepackage{graphicx}
\usepackage{amsmath}
\usepackage{booktabs}
\usepackage[hidelinks]{hyperref}
\usepackage[font=small,labelfont=bf]{caption}

\graphicspath{
  {../outputs/models/rollout_model_comparison/}
  {outputs/models/rollout_model_comparison/}
}

\title{Tangential Phase Error as a Precursor to Catastrophic Channel Escape}
\author{}
\date{}

\begin{document}
\maketitle

\section{Motivation}

Some rollout predictions eventually leave the physically admissible
microfluidic channel. The purpose of this analysis is to test whether these
catastrophic channel escapes are associated with accumulated longitudinal
phase error along the channel trajectory.

Catastrophic escape is defined using the existing channel-admissibility
criterion:
\[
  \mathrm{outside\_fraction} \geq 0.95
\]
at any evaluable rollout step. Each rollout trajectory is therefore grouped as
either escaping or non-escaping. This grouping is based on the predicted
trajectory over the rollout horizon, while the signed phase error is measured
relative to the corresponding true trajectory.

\section{Signed Tangential Error}

For droplet instance \(i\) at rollout step \(k\), the signed tangential
position error is
\[
e_{\mathrm{tan},i}(k)
=
\left(
\hat{\mathbf{x}}_i(k)-\mathbf{x}_i(k)
\right)
\cdot
\mathbf{t}_i(k),
\]
where \(\hat{\mathbf{x}}_i(k)\) is the predicted droplet position,
\(\mathbf{x}_i(k)\) is the true position, and \(\mathbf{t}_i(k)\) is the
consistently oriented local channel tangent evaluated at the true position.

The implementation defines the signed error as predicted position minus true
position, dotted with the true-position tangent. Therefore, positive
\(e_{\mathrm{tan}}\) means the predicted droplet is ahead of the true droplet
along the oriented channel tangent, while negative \(e_{\mathrm{tan}}\) means
the predicted droplet lags behind the true droplet.

At each rollout step, the plotted statistic is the median signed tangential
error over all valid rollout instances in each trajectory class. Because the
validation set contains overlapping rollout windows, these summaries should be
interpreted as statistics over valid rollout instances, not necessarily over
unique physical track IDs.

\section{Result}

\begin{figure}[t]
  \centering
  \includegraphics[width=\linewidth]{phase_direction_stepwise_signed_tangential_error.png}
  \caption{
  Stepwise signed tangential error for catastrophic escaping and non-escaping
  rollout trajectories. Each panel corresponds to one model class:
  \texttt{history\_20}, \texttt{markovian}, and \texttt{geometry\_aware}.
  Curves show the median signed tangential error at each rollout step, and the
  shaded bands show the interquartile range from the 25th to 75th percentiles.
  Negative values indicate predictions lagging behind the true droplet along
  the true-position channel tangent; positive values indicate predictions
  ahead of the true droplet.
  }
  \label{fig:phase-error}
\end{figure}

Figure~\ref{fig:phase-error} shows a clear separation between escaping and
non-escaping rollout trajectories. Non-escaping trajectories remain relatively
close to zero signed tangential error, while escaping trajectories develop a
substantially larger systematic signed tangential error over rollout time.
This separation is visible across all three model classes.

The result indicates that catastrophic channel escapes are not merely isolated
geometric violations at the final moment of failure. They are strongly
associated with an earlier and progressively accumulating longitudinal phase
mismatch between predicted and true droplet motion. This does not prove that
phase error is the sole cause of escape, but it is consistent with phase error
acting as a precursor that likely contributes to later geometric failure.

\section{Physical Interpretation}

A plausible physical mechanism is as follows. First, the model inaccurately
predicts the amount or duration of local slowing. Second, the predicted droplet
becomes out of phase with the true droplet along the channel. Third, subsequent
predicted motion becomes inappropriate for the droplet's current predicted
spatial location. Finally, the trajectory can cross the channel wall and
undergo catastrophic escape.

This phase error can originate from incorrectly learned junction dynamics. At
a junction, the real droplet may slow or stall because of hydrodynamic
interactions and local flow redistribution. If the model does not capture the
correct magnitude or duration of this slowdown, it accumulates longitudinal
phase error. The droplet may then enter subsequent curvature or branch
dynamics from the wrong spatial phase, increasing the risk of wall escape.

Phase mismatch may occur in either direction: the predicted droplet may lag
behind the truth, or it may move ahead of the truth. Both cases can make later
predicted motion inappropriate for the predicted location, particularly near
junctions, bends, and crowded interaction regions.

\section{Implications}

Geometry regularization alone cannot fully correct this failure mode if it
only penalizes the eventual geometric violation. Such a penalty can discourage
wall crossing, but it does not necessarily restore the missing junction or
interaction dynamics that created the accumulated phase error.

The finding motivates a more physics-informed model state. A natural next
step is to include a position-dependent background fluid-flow representation,
with later extensions potentially incorporating droplet morphology or
deformation descriptors. This should be viewed as motivation for future work,
not as a demonstrated solution in the present analysis.

\section{Limitations}

This analysis is associative: trajectories are grouped by whether they
eventually escape, and the observed separation does not establish a unique
causal mechanism. The analysis also does not establish the exact source of the
predicted velocity direction, nor does it prove that phase error is the only
mechanism of catastrophic escape. Finally, because overlapping rollout windows
are included, the summary statistics may reflect repeated rollout instances
from the same physical droplet track.

\end{document}
